\documentclass[10pt,a4paper]{amsart}
\usepackage{amsmath,amssymb,amsthm,graphicx}
% Retained preamble for the Stacks Project chapter snapshot in this directory.
% The chapter's own preamble.tex is a copy of the shared Stacks Project preamble; the block
% below is sliced MECHANICALLY, line for line, from it. Nothing was added, removed or
% redefined; the wrapper class, the xr-hyper/hyperref wiring and the externaldocument list
% are omitted because the assembler supplies the document class.
%
% lines taken from preamble.tex: 3, 5, 9, 11, 12, 156, 157, 158, ... (51 lines total)
%
% package coverage restored from the chapter's own preamble (2026-10-13): xy, xy-2cell

\usepackage[all]{xy}
\xyoption{2cell}
\UseAllTwocells

\newenvironment{reference}{\comment}{\endcomment}
\newenvironment{slogan}{\comment}{\endcomment}
\newenvironment{history}{\comment}{\endcomment}
\theoremstyle{plain}
\newtheorem{theorem}[subsection]{Theorem}
\newtheorem{proposition}[subsection]{Proposition}
\newtheorem{lemma}[subsection]{Lemma}
\theoremstyle{definition}
\newtheorem{definition}[subsection]{Definition}
\newtheorem{example}[subsection]{Example}
\newtheorem{exercise}[subsection]{Exercise}
\newtheorem{situation}[subsection]{Situation}
\theoremstyle{remark}
\newtheorem{remark}[subsection]{Remark}
\newtheorem{remarks}[subsection]{Remarks}
\numberwithin{equation}{subsection}
\def\lim{\mathop{\mathrm{lim}}\nolimits}
\def\colim{\mathop{\mathrm{colim}}\nolimits}
\def\Spec{\mathop{\mathrm{Spec}}}
\def\Hom{\mathop{\mathrm{Hom}}\nolimits}
\def\Ext{\mathop{\mathrm{Ext}}\nolimits}
\def\SheafHom{\mathop{\mathcal{H}\!\mathit{om}}\nolimits}
\def\SheafExt{\mathop{\mathcal{E}\!\mathit{xt}}\nolimits}
\def\Sch{\mathit{Sch}}
\def\Mor{\mathop{\mathrm{Mor}}\nolimits}
\def\Ob{\mathop{\mathrm{Ob}}\nolimits}
\def\Sh{\mathop{\mathit{Sh}}\nolimits}
\def\NL{\mathop{N\!L}\nolimits}
\def\CH{\mathop{\mathrm{CH}}\nolimits}
\def\proetale{{pro\text{-}\acute{e}tale}}
\def\etale{{\acute{e}tale}}
\def\QCoh{\mathit{QCoh}}
\def\Ker{\mathop{\mathrm{Ker}}}
\def\Im{\mathop{\mathrm{Im}}}
\def\Coker{\mathop{\mathrm{Coker}}}
\def\Coim{\mathop{\mathrm{Coim}}}
\DeclareMathSymbol{\boxtimes}{\mathbin}{AMSa}{"02}
\def\QCohstack{\mathcal{QC}\!\mathit{oh}}
\def\Cohstack{\mathcal{C}\!\mathit{oh}}
\def\Spacesstack{\mathcal{S}\!\mathit{paces}}
\def\Quotfunctor{\mathrm{Quot}}
\def\Hilbfunctor{\mathrm{Hilb}}
\def\Curvesstack{\mathcal{C}\!\mathit{urves}}
\def\Polarizedstack{\mathcal{P}\!\mathit{olarized}}
\def\Complexesstack{\mathcal{C}\!\mathit{omplexes}}
\def\Pic{\mathop{\mathrm{Pic}}\nolimits}
\def\Picardstack{\mathcal{P}\!\mathit{ic}}
\def\Picardfunctor{\mathrm{Pic}}
\def\Deformationcategory{\mathcal{D}\!\mathit{ef}}
\usepackage{cleveref}
\crefname{originallemma}{original lemma}{original lemmas}
\Crefname{originallemma}{Original lemma}{Original lemmas}
\makeatletter
\newlabel{lemma-flat-base-change-cotangent-complex}{{08QQ}{source}{}{sourcepointer-1}{}}
\newlabel{lemma-flat-base-change-cotangent-complex@cref}{{[originallemma][0][]08QQ}{[1][0][]source}{}{}{}}
\makeatother
\crefname{originallemma}{original lemma}{original lemmas}
\Crefname{originallemma}{Original lemma}{Original lemmas}
\makeatletter
\newlabel{lemma-special-case-triangle}{{08SA}{source}{}{sourcepointer-2}{}}
\newlabel{lemma-special-case-triangle@cref}{{[originallemma][0][]08SA}{[1][0][]source}{}{}{}}
\makeatother
\title{The cotangent complex of a product of algebras is the sum of the cotangent complexes}
\author{The Stacks Project Authors}
\date{stacks-project:cotangent.tex:local-snapshot-20261012; GNU FDL 1.2 or later}
\begin{document}
\maketitle
\setcounter{subsection}{1}
\noindent\textit{Editorial note.} All delivered mathematical text is the author's own, copied from the pinned Stacks Project chapter file cotangent.tex; no statement or proof step was rewritten, reordered or removed.\par\medskip
\noindent\textit{Editorial note.} The author explicitly chooses the direct elementary proof over the deduction from the fundamental triangle; both the choice and the argument are inside the counted span.\par\medskip
\noindent\textit{Editorial note.} The two references of the counted spans are recorded as non-load-bearing navigation pointers; the argument that consumes them stays inside the counted proof span.\par\medskip
\noindent\textit{Editorial note.} Preamble: the retained preamble of this chapter (sources/stacks-cotangent/retained-preamble.tex, lines 1-63), containing the Xy-pic setup, the comment environments, the theorem declarations and the macros including \textbackslash{}NL.\par\medskip
\noindent\textit{Allowed standard prerequisites (editorial).}\begin{itemize}
\item The special case of the fundamental triangle for a polynomial extension. Lemma lemma-special-case-triangle, used at proof lines 1288-1302. Produces the three distinguished triangles used in the reduction.
\item Flat base change for the cotangent complex. Lemma lemma-flat-base-change-cotangent-complex, used at proof lines 1304-1310. Identifies the cotangent complexes after inverting x.
\item The factorization A -> A[x] -> B x C with x mapping to (1, 0). used at proof lines 1288-1289. Supplies the polynomial extension whose localizations separate B and C.
\item Vanishing of C tensor\_\{A[x]\} A[x, x\textasciicircum{}\{-1\}] and the identification of the localization of B x C. used at proof lines 1305-1307. Make the two localizations reduce to the known cases.
\end{itemize}\medskip
\ifdefined\hypertarget\hypertarget{sourcepointer-1}{}\fi
\noindent\textit{Source reference.} Original lemma 08QQ: focused-seeds/AGP/sources/stacks-snapshot-20261011/tags.txt, source lines 11068--11068. Navigation pointer into the Stacks Project; the standard fact it supplies is applied inside the counted proof span and the argument resting on it stays there.\par\medskip
\ifdefined\hypertarget\hypertarget{sourcepointer-2}{}\fi
\noindent\textit{Source reference.} Original lemma 08SA: focused-seeds/AGP/sources/stacks-snapshot-20261011/tags.txt, source lines 11123--11123. Navigation pointer into the Stacks Project; the standard fact it supplies is applied inside the counted proof span and the argument resting on it stays there.\par\medskip
\begin{lemma}
\label{lemma-cotangent-complex-product}
Let $A \to B$ and $A \to C$ be ring maps.
Then the map $L_{B \times C/A} \to L_{B/A} \oplus L_{C/A}$ is
an isomorphism in $D(B \times C)$.
\end{lemma}

\begin{proof}
Although this lemma can be deduced from the fundamental triangle
we will give a direct and elementary proof of this now.
Factor the ring map $A \to B \times C$ as $A \to A[x] \to B \times C$
where $x \mapsto (1, 0)$. By Lemma \ref{lemma-special-case-triangle}
we have a distinguished triangle
$$
L_{A[x]/A} \otimes_{A[x]}^\mathbf{L} (B \times C) \to L_{B \times C/A} \to
L_{B \times C/A[x]} \to L_{A[x]/A} \otimes_{A[x]}^\mathbf{L} (B \times C)[1]
$$
in $D(B \times C)$. Similarly we have the distinguished triangles
$$
\begin{matrix}
L_{A[x]/A} \otimes_{A[x]}^\mathbf{L} B \to L_{B/A} \to L_{B/A[x]}
\to L_{A[x]/A} \otimes_{A[x]}^\mathbf{L} B[1] \\
L_{A[x]/A} \otimes_{A[x]}^\mathbf{L} C \to L_{C/A} \to L_{C/A[x]}
\to L_{A[x]/A} \otimes_{A[x]}^\mathbf{L} C[1]
\end{matrix}
$$
Thus it suffices to prove the result for $B \times C$ over $A[x]$.
Note that $A[x] \to A[x, x^{-1}]$ is flat, that
$(B \times C) \otimes_{A[x]} A[x, x^{-1}] = B \otimes_{A[x]} A[x, x^{-1}]$,
and that $C \otimes_{A[x]} A[x, x^{-1}] = 0$.
Thus by base change (Lemma \ref{lemma-flat-base-change-cotangent-complex})
the map $L_{B \times C/A[x]} \to L_{B/A[x]} \oplus L_{C/A[x]}$
becomes an isomorphism after inverting $x$.
In the same way one shows that the map becomes an isomorphism after
inverting $x - 1$. This proves the lemma.
\end{proof}

\end{document}
